
\documentclass[../../../physics-standard_questions_all.tex]{subfiles}


\begin{document}

$xy$平面上の電位分布を測定し，2\,V間隔で等電位線を描いたところ，
図のようになった．
等電位線の横に添えた数値は，各等電位線上での電位の値である．

\begin{enumerate}
    
    \item 陽子（電荷$e=1.6 \times 10^{-19}$\,C）を，
    クーロン力に逆らってゆっくりと運ぶとする．
    点Pから点Fへと運ぶ際に必要な仕事$W_\text{PF}$，
    点Pから点Gへと運ぶ際に必要な仕事$W_\text{PG}$，
    点Pから点Hへと運ぶ際に必要な仕事$W_\text{PH}$を，それぞれ求めよ．

    \item 電子（電荷$-e=-1.6 \times 10^{-19}$\,C）を，
    クーロン力に逆らってゆっくりと運ぶとする．
    点Pから点Fへと運ぶ際に必要な仕事$W^*_\text{PF}$，
    点Pから点Gへと運ぶ際に必要な仕事$W^*_\text{PG}$，
    点Pから点Hへと運ぶ際に必要な仕事$W^*_\text{PH}$を，それぞれ求めよ．

    \item 点Pを通る電気力線の概略を描け．
    
    \item 点P付近に図示してある等電位線の間隔は大体0.53\,m程度である．
    このことから点Pにおける電場の強さ$E_\text{P}$の近似値を求めよ．

\end{enumerate}

\vspace{0.2\baselineskip}
{\centering
    \includegraphics{\subfix{fig/fig_em_ef_04/fig_em_ef_04_03.pdf}}
\par}
\vspace{0.2\baselineskip}

\end{document}



